\documentclass[10pt,a4paper]{amsart}
\usepackage{amsmath,amssymb,amsthm,graphicx}
\usepackage{verbatim}
\newenvironment{reference}{\comment}{\endcomment}
%\newenvironment{reference}{}{}
\newenvironment{slogan}{\comment}{\endcomment}
\newenvironment{history}{\comment}{\endcomment}
% For commutative diagrams we use Xy-pic
\usepackage[all]{xy}

% We use 2cell for 2-commutative diagrams.
\xyoption{2cell}
% Theorem environments.
%
\theoremstyle{plain}
\newtheorem{theorem}[subsection]{Theorem}
\newtheorem{proposition}[subsection]{Proposition}
\newtheorem{lemma}[subsection]{Lemma}

\theoremstyle{definition}
\newtheorem{definition}[subsection]{Definition}
\newtheorem{example}[subsection]{Example}
\newtheorem{exercise}[subsection]{Exercise}
\newtheorem{situation}[subsection]{Situation}

\theoremstyle{remark}
\newtheorem{remark}[subsection]{Remark}
\newtheorem{remarks}[subsection]{Remarks}

\numberwithin{equation}{subsection}

% Macros
%
\def\lim{\mathop{\mathrm{lim}}\nolimits}
\def\colim{\mathop{\mathrm{colim}}\nolimits}
\def\Spec{\mathop{\mathrm{Spec}}}
\def\Hom{\mathop{\mathrm{Hom}}\nolimits}
\def\Ext{\mathop{\mathrm{Ext}}\nolimits}
\def\SheafHom{\mathop{\mathcal{H}\!\mathit{om}}\nolimits}
\def\SheafExt{\mathop{\mathcal{E}\!\mathit{xt}}\nolimits}
\def\Sch{\mathit{Sch}}
\def\Mor{\mathop{\mathrm{Mor}}\nolimits}
\def\Ob{\mathop{\mathrm{Ob}}\nolimits}
\def\Sh{\mathop{\mathit{Sh}}\nolimits}
\def\NL{\mathop{N\!L}\nolimits}
\def\CH{\mathop{\mathrm{CH}}\nolimits}
\def\proetale{{pro\text{-}\acute{e}tale}}
\def\etale{{\acute{e}tale}}
\def\QCoh{\mathit{QCoh}}
\def\Ker{\mathop{\mathrm{Ker}}}
\def\Im{\mathop{\mathrm{Im}}}
\def\Coker{\mathop{\mathrm{Coker}}}
\def\Coim{\mathop{\mathrm{Coim}}}

% Boxtimes
%
\DeclareMathSymbol{\boxtimes}{\mathbin}{AMSa}{"02}

%
% Macros for moduli stacks/spaces
%
\def\QCohstack{\mathcal{QC}\!\mathit{oh}}
\def\Cohstack{\mathcal{C}\!\mathit{oh}}
\def\Spacesstack{\mathcal{S}\!\mathit{paces}}
\def\Quotfunctor{\mathrm{Quot}}
\def\Hilbfunctor{\mathrm{Hilb}}
\def\Curvesstack{\mathcal{C}\!\mathit{urves}}
\def\Polarizedstack{\mathcal{P}\!\mathit{olarized}}
\def\Complexesstack{\mathcal{C}\!\mathit{omplexes}}
% \Pic is the operator that assigns to X its picard group, usage \Pic(X)
% \Picardstack_{X/B} denotes the Picard stack of X over B
% \Picardfunctor_{X/B} denotes the Picard functor of X over B
\def\Pic{\mathop{\mathrm{Pic}}\nolimits}
\def\Picardstack{\mathcal{P}\!\mathit{ic}}
\def\Picardfunctor{\mathrm{Pic}}
\def\Deformationcategory{\mathcal{D}\!\mathit{ef}}
\usepackage{cleveref}
\crefname{originallemma}{original lemma}{original lemmas}
\Crefname{originallemma}{Original lemma}{Original lemmas}
\makeatletter
\newlabel{lemma-abelian}{{lemma-abelian}{source}{}{sourcepointer-1}{}}
\newlabel{lemma-abelian@cref}{{[originallemma][0][]lemma-abelian}{[1][0][]source}{}{}{}}
\makeatother
\crefname{originallemma}{original lemma}{original lemmas}
\Crefname{originallemma}{Original lemma}{Original lemmas}
\makeatletter
\newlabel{lemma-stalk-module-differentials}{{lemma-stalk-module-differentials}{source}{}{sourcepointer-2}{}}
\newlabel{lemma-stalk-module-differentials@cref}{{[originallemma][0][]lemma-stalk-module-differentials}{[1][0][]source}{}{}{}}
\makeatother
\crefname{originalreference}{original reference}{original references}
\Crefname{originalreference}{Original reference}{Original references}
\makeatletter
\newlabel{algebra-lemma-NL-homotopy}{{algebra-lemma-NL-homotopy}{source}{}{sourcepointer-3}{}}
\newlabel{algebra-lemma-NL-homotopy@cref}{{[originalreference][0][]algebra-lemma-NL-homotopy}{[1][0][]source}{}{}{}}
\makeatother
\crefname{originalreference}{original reference}{original references}
\Crefname{originalreference}{Original reference}{Original references}
\makeatletter
\newlabel{algebra-section-netherlander}{{algebra-section-netherlander}{source}{}{sourcepointer-4}{}}
\newlabel{algebra-section-netherlander@cref}{{[originalreference][0][]algebra-section-netherlander}{[1][0][]source}{}{}{}}
\makeatother
\title{The naive cotangent complex does not depend on the chosen presentation}
\author{The Stacks Project Authors}
\date{stacks-project:modules.tex:local-snapshot-20261012; GNU FDL 1.2 or later}
\begin{document}
\maketitle
\noindent\textit{Editorial note.} All delivered mathematical text is the author's own, copied from the pinned Stacks Project chapter source sources/stacks-modules/pinned-original-source.tex; no statement or proof step was rewritten, reordered or removed.\par\medskip
\noindent\textit{Editorial note.} Preamble provenance: the Stacks Project keeps one single shared preamble for every chapter; this card takes it directly from the staged sources/stacks-modules/preamble.tex as three verbatim spans (lines 8-12, 14-18 and 154-217), so no derived file was written and nothing was added or redefined. The wrapper class and the base packages are supplied by the assembler.\par\medskip
\noindent\textit{Editorial note.} Bibliography provenance: neither the counted statement nor its proof contains a citation (this lane rescans the whole span before writing the card), so the card carries no bibliography and no entry was reconstructed.\par\medskip
\noindent\textit{Editorial note.} Selection provenance: this card was opened under the file-based control contract (CONTROL.md), with the mechanical pre-card gate run on its statement and proof interval before the card was written.\par\medskip
\noindent\textit{Allowed standard prerequisites (editorial).}\begin{itemize}
\item The complex attached to a presentation of a sheaf of algebras as a quotient of a polynomial algebra, and the complex defined by the identity presentation. The set-up paragraph is reproduced verbatim as a context span of this card (source lines 5922-5939). Fixes the two objects that the statement compares.
\item The compatibility of the formation of differentials with taking stalks, and the description of the naive cotangent complex of a presentation over a ring. Invoked inside the counted proof at source lines 5966-5972 as pointers to a lemma of this chapter and to a section of another chapter. Reduces the required quasi-isomorphism to the corresponding statement in commutative algebra.
\end{itemize}\medskip
\ifdefined\hypertarget\hypertarget{sourcepointer-1}{}\fi
\noindent\textit{Source reference.} Original lemma lemma-abelian: sources/stacks-modules/pinned-original-source.tex, source lines 439--439. Navigation pointer only: the counted text refers to the Stacks item lemma-abelian, which belongs to another chapter of the project and is not part of this card. The pointer carries the target's own label as its number because this chapter's staged snapshot does not contain the other chapter's tag table; its content is not reproduced here.\par\medskip
\ifdefined\hypertarget\hypertarget{sourcepointer-2}{}\fi
\noindent\textit{Source reference.} Original lemma lemma-stalk-module-differentials: sources/stacks-modules/pinned-original-source.tex, source lines 5197--5197. Navigation pointer only: the counted text refers to the Stacks item lemma-stalk-module-differentials, which belongs to another chapter of the project and is not part of this card. The pointer carries the target's own label as its number because this chapter's staged snapshot does not contain the other chapter's tag table; its content is not reproduced here.\par\medskip
\ifdefined\hypertarget\hypertarget{sourcepointer-3}{}\fi
\noindent\textit{Source reference.} Original reference algebra-lemma-NL-homotopy: sources/stacks-modules/pinned-original-source.tex, source lines 5974--5974. Navigation pointer only: the counted text refers to the Stacks item algebra-lemma-NL-homotopy, which belongs to another chapter of the project and is not part of this card. The pointer carries the target's own label as its number because this chapter's staged snapshot does not contain the other chapter's tag table; its content is not reproduced here.\par\medskip
\ifdefined\hypertarget\hypertarget{sourcepointer-4}{}\fi
\noindent\textit{Source reference.} Original reference algebra-section-netherlander: sources/stacks-modules/pinned-original-source.tex, source lines 5827--5827. Navigation pointer only: the counted text refers to the Stacks item algebra-section-netherlander, which belongs to another chapter of the project and is not part of this card. The pointer carries the target's own label as its number because this chapter's staged snapshot does not contain the other chapter's tag table; its content is not reproduced here.\par\medskip
\begin{lemma}
\label{lemma-NL-up-to-qis}
In the situation above there is a canonical isomorphism
$\NL(\alpha) = \NL_{\mathcal{B}/\mathcal{A}}$ in $D(\mathcal{B})$.
\end{lemma}

\medskip\noindent
We can choose a different presentation of $\mathcal{B}$ as a quotient of a
polynomial algebra over $\mathcal{A}$ and still obtain the same object
of $D(\mathcal{B})$. To explain this, suppose that $\mathcal{E}$ is
a sheaves of sets on $X$ and $\alpha : \mathcal{E} \to \mathcal{B}$
a map of sheaves of sets. Then we obtain an $\mathcal{A}$-algebra
homomorphism $\mathcal{A}[\mathcal{E}] \to \mathcal{B}$. If this map
is surjective, i.e., if $\alpha(\mathcal{E})$ generates $\mathcal{B}$
as an $\mathcal{A}$-algebra, then we set
$$
\NL(\alpha) = \left(
\mathcal{J}/\mathcal{J}^2
\longrightarrow
\Omega_{\mathcal{A}[\mathcal{E}]/\mathcal{A}}
\otimes_{\mathcal{A}[\mathcal{E}]} \mathcal{B}\right)
$$
where $\mathcal{J} \subset \mathcal{A}[\mathcal{E}]$ is the kernel
of the surjection $\mathcal{A}[\mathcal{E}] \to \mathcal{B}$.

\begin{proof}
Observe that $\NL_{\mathcal{B}/\mathcal{A}} = \NL(\text{id}_\mathcal{B})$.
Thus it suffices to show that given two maps
$\alpha_i : \mathcal{E}_i \to \mathcal{B}$ as above, there is a
canonical quasi-isomorphism $\NL(\alpha_1) = \NL(\alpha_2)$ in $D(\mathcal{B})$.
To see this set $\mathcal{E} = \mathcal{E}_1 \amalg \mathcal{E}_2$ and
$\alpha = \alpha_1 \amalg \alpha_2 : \mathcal{E} \to \mathcal{B}$.
Set
$\mathcal{J}_i = \Ker(\mathcal{A}[\mathcal{E}_i] \to \mathcal{B})$
and
$\mathcal{J} = \Ker(\mathcal{A}[\mathcal{E}] \to \mathcal{B})$.
We obtain maps $\mathcal{A}[\mathcal{E}_i] \to \mathcal{A}[\mathcal{E}]$
which send $\mathcal{J}_i$ into $\mathcal{J}$.
Thus we obtain canonical maps of complexes
$$
\NL(\alpha_i) \longrightarrow \NL(\alpha)
$$
and it suffices to show these maps are quasi-isomorphism. To see this
it suffices to check on stalks (Lemma \ref{lemma-abelian}). If $x \in X$
then the stalk of $\NL(\alpha)$ is the complex $\NL(\alpha_x)$ of
Algebra, Section \ref{algebra-section-netherlander}
associated to the presentation $\mathcal{A}_x[\mathcal{E}_x] \to \mathcal{B}_x$
coming from the map $\alpha_x : \mathcal{E}_x \to \mathcal{B}_x$.
(Some details omitted; use Lemma \ref{lemma-stalk-module-differentials}
to see compatibility of forming differentials and taking stalks.)
We conclude the result holds by
Algebra, Lemma \ref{algebra-lemma-NL-homotopy}.
\end{proof}

\end{document}
